%*******************************************************
% Abstract
%*******************************************************
\pdfbookmark[0]{Zusammenfassung}{Zusammenfassung}
\chapter*{Zusammenfassung}

Die Arbeit untersucht die Multiclass-Klassifikation von sieben Dry-Bean-Sorten mit acht morphologischen Merkmalen aus dem UCI Dry Bean Dataset. Verglichen werden Logistic Regression, Random Forest, HistGradientBoosting und ein Multi-Layer-Perzeptron. Die Daten werden stratifiziert in Trainings- und Testdaten aufgeteilt. Die Bewertung erfolgt mit Accuracy, Macro-F1 und Konfusionsmatrizen, da die Klassen ungleich verteilt sind.

Alle Modelle erreichen hohe und nahe beieinanderliegende Testwerte. Das MLP erzielt mit einer Accuracy von 0{,}923 und einem Macro-F1 von 0{,}934 das beste Ergebnis. HistGradientBoosting erreicht 0{,}920 und 0{,}932, Random Forest 0{,}918 und 0{,}931, die logistische Regression 0{,}917 und 0{,}931. Die Fehler konzentrieren sich vor allem auf morphologisch ähnliche Sorten, insbesondere DERMASON und SIRA sowie BARBUNYA und CALI.

Zur Erklärung werden Permutation Feature Importance, SHAP und LIME auf alle vier finalen Modelle angewandt. Die globalen Vergleiche zeigen über die Modelle hinweg eine wiederkehrende Bedeutung von \texttt{Area} und \texttt{ShapeFactor1}. Lokale Erklärungen behandeln eine korrekt klassifizierte DERMASON-Instanz sowie DERMASON-als-SIRA-Fehlfälle. Die Erklärungen passen zu der Annahme, dass die Modelle Größen- und Forminformationen der Bohnen nutzen. Wegen korrelierter Merkmale, der begrenzten lokalen Beispiele und eines einzelnen Testsplit sind die Rangfolgen jedoch vorsichtig zu interpretieren. Die XAI-Ergebnisse beschreiben Modellverhalten und keine kausalen Zusammenhänge.


%Kurze Zusammenfassung des Inhaltes in deutscher Sprache, der Umfang beträgt zwischen einer halben und einer ganzen DIN A4-Seite.

%Orientieren Sie sich bei der Aufteilung bzw. dem Inhalt Ihrer Zusammenfassung an Kent Becks Artikel: \url{http://plg.uwaterloo.ca/~migod/research/beckOOPSLA.html}.
